%%%%%%%%%%%%%%%%%%%%%%%%%%%%%%%%%%%%%%%%%%%%%%%%%%%%%%%%%%%%%%%%%%%%%%%%%%%%%%%
% CHAPTER 10: WELFARE AND THE FEPSDE FRAMEWORK
%%%%%%%%%%%%%%%%%%%%%%%%%%%%%%%%%%%%%%%%%%%%%%%%%%%%%%%%%%%%%%%%%%%%%%%%%%%%%%%
% Version: 2.0 (January 2026)
%
% Purpose: Introduces the 144-component utility structure (INU/KNU/IDN × FEPSDE
%          × Time × Valence) answering WHAT constitutes utility and WHO has it.
%
% Primary Appendix: C - CORE-WHAT (The FEPSDE Matrix)
% Secondary Appendices: AAA (CORE-WHO), B (CORE-HOW), AU (BCM Axioms), AJ (Social)
%
% Prerequisites: Ch. 5 (Complementarity), Ch. 8 (Mathematical Foundations)
% Leads to: Ch. 11 (Awareness), Ch. 12 (Willingness)
%
% Chapter Type: A (CORE Chapter - answers WHAT and WHO)
%%%%%%%%%%%%%%%%%%%%%%%%%%%%%%%%%%%%%%%%%%%%%%%%%%%%%%%%%%%%%%%%%%%%%%%%%%%%%%%

\chapter{Welfare and the FEPSDE Framework: The 144-Component Utility Structure}
\label{ch:welfare-fepsde}

%%%%%%%%%%%%%%%%%%%%%%%%%%%%%%%%%%%%%%%%%%%%%%%%%%%%%%%%%%%%%%%%%%%%%%%%%%%%%%%
% CHAPTER SCOPE BOX
%%%%%%%%%%%%%%%%%%%%%%%%%%%%%%%%%%%%%%%%%%%%%%%%%%%%%%%%%%%%%%%%%%%%%%%%%%%%%%%
\begin{tcolorbox}[
    colback=purple!5!white,
    colframe=purple!75!black,
    title={\textbf{Chapter Scope: Ziel / In-Scope / Out-of-Scope}},
    fonttitle=\bfseries
]

\textbf{Ziel (Objective):}
Establish the formal foundation for inu/knu/idn, 144 components, fepsde within the 10C CORE framework.

\vspace{0.3cm}
\textbf{In-Scope:}
\begin{itemize}[nosep]
    \item Core definitions and axioms for inu/knu/idn
    \item Integration with other 10C dimensions
    \item Formal mathematical foundations
    \item Key theorems and their implications
\end{itemize}

\vspace{0.3cm}
\textbf{Out-of-Scope (delegiert an Appendices):}
\begin{itemize}[nosep]
    \item Detailed proofs and derivations $\rightarrow$ CORE Appendix
    \item Empirical calibration $\rightarrow$ Appendix BBB
    \item Domain-specific applications $\rightarrow$ Application chapters
\end{itemize}

\end{tcolorbox}



%%%%%%%%%%%%%%%%%%%%%%%%%%%%%%%%%%%%%%%%%%%%%%%%%%%%%%%%%%%%%%%%%%%%%%%%%%%%%%%
% QUICK REFERENCE BOX
%%%%%%%%%%%%%%%%%%%%%%%%%%%%%%%%%%%%%%%%%%%%%%%%%%%%%%%%%%%%%%%%%%%%%%%%%%%%%%%
\begin{tcolorbox}[
    colback=gray!5!white,
    colframe=gray!75!black,
    title={\textbf{Begriffe in diesem Kapitel}},
    fonttitle=\bfseries
]
\small
\begin{tabular}{@{}ll@{}}
\textbf{Key Concept} & \textbf{Definition} \\
\midrule
INU/KNU/IDN & See Section \ref{sec:ch10-intro} \\
\end{tabular}
\end{tcolorbox}



%%%%%%%%%%%%%%%%%%%%%%%%%%%%%%%%%%%%%%%%%%%%%%%%%%%%%%%%%%%%%%%%%%%%%%%%%%%%%%%
% INTUITION BOX
%%%%%%%%%%%%%%%%%%%%%%%%%%%%%%%%%%%%%%%%%%%%%%%%%%%%%%%%%%%%%%%%%%%%%%%%%%%%%%%
\begin{tcolorbox}[
    colback=blue!5!white,
    colframe=blue!75!black,
    title={\textbf{Intuition: A Motivating Example}},
    fonttitle=\bfseries
]
Consider \textbf{Anna}, a professional facing a decision that involves inu/knu/idn, 144 components, fepsde.

She initially evaluates the choice using standard criteria, but something feels incomplete.
\textbf{Thomas}, her colleague, suggests considering additional dimensions beyond the obvious.

This chapter explores what Anna discovers when she applies the EBF framework---how
inu/knu/idn, 144 components, fepsde interact with broader welfare considerations.

\medskip
\textit{By the end of this chapter, you will understand why Anna's initial analysis
was incomplete and how the EBF approach provides a more comprehensive view.}
\end{tcolorbox}



%==============================================================================
% DIE 9 FRAGEN NAVIGATION BOX
%==============================================================================
\BCMQuestionNav{10}

%==============================================================================
% 10C CORE CONNECTION BOX
%==============================================================================
\begin{tcolorbox}[
    colback=red!5!white,
    colframe=red!75!black,
    title={\textbf{10C CORE Connection: WHAT (Question 3) + WHO (Question 1)}},
    fonttitle=\bfseries
]
\textbf{The Fundamental Questions:}
\begin{quote}
\textbf{WHAT:} What dimensions constitute utility? $\rightarrow$ FEPSDE (6 dimensions)\\
\textbf{WHO:} At which level is utility evaluated? $\rightarrow$ $L \in \{0, 1, ..., \infty\}$
\end{quote}

\medskip
\begin{tabular}{ll}
\textbf{CORE Appendices:} & C (CORE-WHAT) + AAA (CORE-WHO) \\
\textbf{Output:} & 144-component utility structure: $3 \times 6 \times 4 \times 2$ \\
\textbf{Key Insight:} & Coherence ($K$) without Quality ($Q$) is incomplete \\
\textbf{Novel Contribution:} & INU/KNU/IDN categories with level-indexed KNU$^L$
\end{tabular}

\medskip
\textbf{The 10C in Chapter 10:}
\begin{center}
\small
\begin{tabular}{@{}clll@{}}
\toprule
\textbf{Section} & \textbf{CORE} & \textbf{Question} & \textbf{Output} \\
\midrule
10.1--10.3 & WHAT (C) & What is utility? & FEPSDE, 144 components \\
10.4 & HOW (B) & Inter-category $\gamma$ & Complementarities \\
10.5 & WHEN (V) & Context modulation & $\Gamma$-matrix \\
10.6 & WHO (AAA) & Level aggregation & $\alpha^L(\Psi)$ \\
10.7 & STAGE (AW) & BCJ integration & Stage-dependent weights \\
\bottomrule
\end{tabular}
\end{center}

\textit{This chapter answers the two most fundamental questions: What do people value? (FEPSDE) and For whom? (Level hierarchy). Without these, complementarity (HOW) and context (WHEN) have nothing to operate on.}
\end{tcolorbox}

%==============================================================================
% APPENDIX REFERENCES BOX
%==============================================================================
\begin{tcolorbox}[
    colback=yellow!10!white,
    colframe=orange!75!black,
    title={\textbf{Appendix References for This Chapter}},
    fonttitle=\bfseries
]
\begin{tabular}{lll}
\textbf{Appendix} & \textbf{Content} & \textbf{Section} \\
\midrule
C (CORE-WHAT) & Complete 144-component structure, Meta-Axiom U-M0 & \ref{sec:fepsde-complete} \\
AAA (CORE-WHO) & Aggregation levels ($L = 0...\infty$) & \ref{sec:fepsde-levels} \\
B (CORE-HOW) & Inter-category complementarities & \ref{sec:fepsde-complementarity} \\
AU (BCM Axioms) & Complete formal axiomatization & All sections \\
AJ (Social Preferences) & INU/KNU/IDN foundations & \ref{sec:fepsde-categories} \\
G (REF-GLOSSARY) & Symbol definitions & All sections
\end{tabular}

\medskip
\textbf{Reading Guide:} This chapter introduces the FEPSDE dimensions conceptually;
Appendix WAT provides the complete 144-component formal specification.
\end{tcolorbox}

%==============================================================================
% CROSS-REFERENCE: EIT-EMP / EMERGE ALGORITHM
%==============================================================================
\begin{tcolorbox}[
    colback=blue!5!white,
    colframe=blue!75!black,
    title={\textbf{Cross-Reference: EMERGE Algorithm (Chapter 11, EIT-EMP)}},
    fonttitle=\bfseries
]
\textbf{WHAT $\rightarrow$ Intervention Design:} Die 144-Komponenten-Utility-Struktur dieses Kapitels definiert den \textbf{Zielraum} für Interventionen. Das EMERGE-Algorithmus (Chapter 11) nutzt diese Struktur zur Berechnung optimaler Interventionsprofile.

\medskip
\textbf{Dimensionsmapping:}
\begin{center}
\begin{tabular}{lll}
\textbf{Chapter 10 (WHAT)} & \textbf{Intervention-Vektor} & \textbf{EMERGE-Nutzung} \\
\midrule
INU (Intrapersonal) & $I_{\text{WHAT},\text{INU}}$ & Selbstbezogene Interventionen \\
KNU (Relational) & $I_{\text{WHAT},\text{KNU}}$ & Soziale Norm-Interventionen \\
IDN (Identity) & $I_{\text{WHAT},\text{IDN}}$ & Identitäts-Interventionen \\
\midrule
Financial (F) & $\vec{I}_F \in [0,1]$ & Monetäre Anreize \\
Emotional (E) & $\vec{I}_E \in [0,1]$ & Affektive Interventionen \\
Physical (P) & $\vec{I}_P \in [0,1]$ & Gesundheits-Nudges \\
Social (S) & $\vec{I}_S \in [0,1]$ & Peer-Effekte, Normen \\
Domain (D) & $\vec{I}_D \in [0,1]$ & Bereichsspezifische Maßnahmen \\
Existential (X) & $\vec{I}_X \in [0,1]$ & Sinn-/Purpose-Interventionen \\
\end{tabular}
\end{center}

\medskip
\textbf{EMERGE-Formel (EIT-EMP-2, EXC-3 Hybrid):}
\[
K^* = K_{\text{base}} \cdot f_B + \delta_C + \delta_S + \delta_\gamma + \delta_L
\]
Nur $f_B$ ist multiplikativ (Veto-Logik: $B=0 \Rightarrow K^*=0$). Der Komplexitäts-Adjustment $\delta_C$ ist additiv und hängt von der Anzahl adressierter FEPSDE-Dimensionen ab:
\[
f_C = \begin{cases}
1.0 & \text{wenn } |\text{FEPSDE}_{\text{target}}| \leq 2 \\
0.85 & \text{wenn } |\text{FEPSDE}_{\text{target}}| = 3-4 \\
0.7 & \text{wenn } |\text{FEPSDE}_{\text{target}}| \geq 5
\end{cases}
\]

\medskip
\textbf{Workflow:} Chapter 10 (WHAT: $u_d$, INU/KNU/IDN) $\rightarrow$ Chapter 11 (EMERGE: $K^*$, $P^*$) $\rightarrow$ Chapters 17--22 (Implementation)

\medskip
\textit{Siehe Chapter 11, Section 11.5 für die vollständige EMERGE-Spezifikation.}
\end{tcolorbox}


%%%%%%%%%%%%%%%%%%%%%%%%%%%%%%%%%%%%%%%%%%%%%%%%%%%%%%%%%%%%%%%%%%%%%%%%%%%%%%%
%                     PART I: THE PROBLEM OF WELFARE
%%%%%%%%%%%%%%%%%%%%%%%%%%%%%%%%%%%%%%%%%%%%%%%%%%%%%%%%%%%%%%%%%%%%%%%%%%%%%%%

\section{Introduction: Why Coherence Is Not Enough}
\label{sec:fepsde-intro}

%------------------------------------------------------------------------------

The coherence index $K$ measures whether a system is in equilibrium.
It does not measure whether the equilibrium is \emph{good}.
A society of mutual exploitation can be perfectly coherent ($K = 1$)
yet deeply impoverished in welfare.

\begin{intuition}[Two Coherent Societies]
Consider two societies, both with $K \approx 1$:

\textbf{Society A: Extractive Coherence}
\begin{itemize}[nosep]
  \item Elites extract resources from masses
  \item Masses expect extraction and don't invest
  \item Institutions enforce extraction
  \item Everyone's beliefs match reality---coherent
  \item But welfare is low for most people
\end{itemize}

\textbf{Society B: Inclusive Coherence}
\begin{itemize}[nosep]
  \item Property rights protect investment
  \item People expect returns and invest
  \item Institutions enforce contracts
  \item Everyone's beliefs match reality---coherent
  \item Welfare is high for most people
\end{itemize}

Both have $K \approx 1$. But B is clearly better than A.
We need a criterion to capture this: the \textbf{quality index $Q$}.
\end{intuition}

%------------------------------------------------------------------------------
% CENTRAL QUESTION BOX
%------------------------------------------------------------------------------
\paragraph{The Central Questions.}
This chapter addresses two of the 9 BCM Questions:

\begin{center}
\fcolorbox{red!75!black}{red!5!white}{\parbox{0.85\textwidth}{
\textbf{Question 1 (WHO):} At which level is utility evaluated?\\[0.5em]
\textbf{Question 3 (WHAT):} What dimensions constitute utility?
}}
\end{center}

\paragraph{Why Not Just Use GDP?}
Gross Domestic Product measures market transactions. It misses:
\begin{itemize}[nosep]
  \item Non-market production (household work, volunteering)
  \item Distribution (GDP can rise while most people suffer)
  \item Health and longevity
  \item Environmental degradation
  \item Social cohesion and trust
  \item Subjective wellbeing
  \item Sustainability (GDP today vs. GDP tomorrow)
\end{itemize}

A framework that claims to capture ``rationality'' must have a richer welfare concept.

%------------------------------------------------------------------------------
% CHAPTER OVERVIEW
%------------------------------------------------------------------------------
\paragraph{Chapter Overview.}
This chapter proceeds as follows:
\begin{itemize}
    \item Section \ref{sec:fepsde-dimensions}: The six FEPSDE dimensions
    \item Section \ref{sec:fepsde-categories}: The three utility categories (INU/KNU/IDN)
    \item Section \ref{sec:fepsde-time}: Time structure and discounting
    \item Section \ref{sec:fepsde-valence}: Gains, pains, and loss aversion
    \item Section \ref{sec:fepsde-complete}: The complete 144-component structure
    \item Section \ref{sec:fepsde-applications}: Country comparisons and applications
    \item Section \ref{sec:fepsde-summary}: Summary and transition
\end{itemize}


%%%%%%%%%%%%%%%%%%%%%%%%%%%%%%%%%%%%%%%%%%%%%%%%%%%%%%%%%%%%%%%%%%%%%%%%%%%%%%%
%                     PART II: THE SIX FEPSDE DIMENSIONS
%%%%%%%%%%%%%%%%%%%%%%%%%%%%%%%%%%%%%%%%%%%%%%%%%%%%%%%%%%%%%%%%%%%%%%%%%%%%%%%

\section{The FEPSDE Dimensions: What Do People Value?}
\label{sec:fepsde-dimensions}

%------------------------------------------------------------------------------

\paragraph{The Selection Criteria.}
The dimensions were selected based on:
\begin{enumerate}[nosep]
  \item \textbf{Empirical distinctness:} Each captures variance not captured by others
  \item \textbf{Universal relevance:} Applies across cultures, levels, time periods
  \item \textbf{Measurability:} Can be operationalized for empirical work
  \item \textbf{Policy relevance:} Corresponds to distinct policy domains
  \item \textbf{Theoretical grounding:} Connects to established research programs
\end{enumerate}

\begin{definition}[FEPSDE Utility Dimensions]
\label{def:fepsde}
\begin{align}
  U^F &: \text{\textbf{Financial} --- material wealth, income, economic security} \\
  U^E &: \text{\textbf{Emotional} --- psychological wellbeing, life satisfaction, meaning} \\
  U^P &: \text{\textbf{Physical} --- health, energy, bodily integrity, longevity} \\
  U^S &: \text{\textbf{Social} --- relationships, belonging, trust, status, community} \\
  U^D &: \text{\textbf{Digital} --- connectivity, access, digital participation, data rights} \\
  U^{Eco} &: \text{\textbf{Ecological} --- environmental quality, sustainability, nature access}
\end{align}
\end{definition}

%------------------------------------------------------------------------------
\subsection{Financial Dimension ($U^F$)}
\label{sec:fepsde-financial}

\begin{definition}[Financial Utility]
Material resources and economic security: income, wealth, consumption possibilities,
economic security, property, and access to financial services.
\end{definition}

\paragraph{Characteristics.}
\begin{itemize}[nosep]
  \item Core of classical economic analysis
  \item Necessary but not sufficient for wellbeing
  \item Subject to diminishing marginal utility
  \item Relative position matters (status concerns)
\end{itemize}

\paragraph{Why It's Not Sufficient Alone.}
The ``Easterlin Paradox'': above a threshold, more income doesn't increase happiness.
Cross-country: GDP per capita correlates with wellbeing only up to ~\$20,000.
Within countries: relative position matters more than absolute level.

\begin{intuition}[High $U^F$, Low Other Dimensions]
A wealthy executive: high income, long hours, no time for family,
chronic stress, sedentary lifestyle, no community involvement.
$U^F$ high; $U^E$, $U^P$, $U^S$ low.
\end{intuition}

\paragraph{Measurement.}
Income (individual, household, equivalized), wealth (assets minus liabilities),
consumption expenditure, economic security indices, Gini coefficient (distribution).

\begin{cross-reference}
See Appendix WAT (CORE-WHAT), Section C.2 for the complete axiomatization of
the Financial dimension including loss aversion parameter $\lambda_F \approx 2$.
\end{cross-reference}

%------------------------------------------------------------------------------
\subsection{Emotional Dimension ($U^E$)}
\label{sec:fepsde-emotional}

\begin{definition}[Emotional Utility]
Psychological wellbeing and subjective experience: life satisfaction,
positive and negative affect, meaning and purpose, autonomy, mental health.
\end{definition}

\paragraph{Theoretical Grounding.}
Connects to:
\begin{itemize}[nosep]
  \item Subjective wellbeing research \citep{diener1999}
  \item Self-determination theory (autonomy, competence, relatedness)
  \item Positive psychology (PERMA: Positive emotion, Engagement, Relationships, Meaning, Achievement)
\end{itemize}

\begin{intuition}[The Nordic Paradox]
Nordic countries: moderate $U^F$ (high taxes reduce disposable income)
but very high $U^E$ (life satisfaction, low anxiety, high trust).
This demonstrates that $U^E$ is empirically distinct from $U^F$.
\end{intuition}

\paragraph{Measurement.}
Life satisfaction scales (Cantril ladder, SWLS), positive/negative affect schedules (PANAS),
depression and anxiety inventories, meaning in life questionnaires,
experience sampling (momentary wellbeing).

%------------------------------------------------------------------------------
\subsection{Physical Dimension ($U^P$)}
\label{sec:fepsde-physical}

\begin{definition}[Physical Utility]
Bodily health and functioning: health status, physical functioning, mobility,
energy, vitality, longevity, freedom from pain, sleep quality.
\end{definition}

\paragraph{Why It's Fundamental.}
Health is a precondition for enjoying other dimensions.
Severe illness reduces capacity to work ($U^F$), enjoy relationships ($U^S$),
and experience positive emotions ($U^E$).

\begin{intuition}[The American Health Paradox]
USA: Highest healthcare spending per capita.
Yet: Lower life expectancy than peer countries, high rates of chronic disease.
High $U^F$ investment in health doesn't automatically produce high $U^P$.
\end{intuition}

\paragraph{Measurement.}
Life expectancy, healthy life years, self-reported health, DALYs, QALYs,
biomarkers (blood pressure, BMI), functional capacity assessments.

%------------------------------------------------------------------------------
\subsection{Social Dimension ($U^S$)}
\label{sec:fepsde-social}

\begin{definition}[Social Utility]
Relationships and social integration: close relationships (family, friends, partners),
social support networks, community belonging, social status, trust, civic participation.
\end{definition}

\paragraph{Theoretical Grounding.}
Connects to:
\begin{itemize}[nosep]
  \item Social capital literature \citep{putnam2000}
  \item Attachment theory
  \item Belongingness as fundamental need
  \item Status and relative position research
\end{itemize}

\paragraph{Why Social Relationships Are Special.}
Strongest predictor of longevity (stronger than smoking, obesity, exercise).
Loneliness is a public health crisis. Social trust underpins economic transactions.

\begin{intuition}[Japan's Social Crisis]
Japan: High $U^F$, high $U^P$ (longevity), but increasingly low $U^S$.
``Hikikomori'' (social withdrawal), declining marriage, ``kodokushi'' (lonely deaths).
Economic success doesn't guarantee social wellbeing.
\end{intuition}

\paragraph{Measurement.}
Social network size and quality, frequency of social contact,
loneliness scales (UCLA Loneliness Scale), trust surveys (World Values Survey),
civic participation rates, social support indices.

%------------------------------------------------------------------------------
\subsection{Digital Dimension ($U^D$)}
\label{sec:fepsde-digital}

\begin{definition}[Digital Utility]
Participation in the digital world: internet access, digital literacy,
access to digital services, data rights and privacy, protection from digital harms,
ability to disconnect (digital wellbeing).
\end{definition}

\paragraph{Why a Separate Dimension?}
Digital participation has become:
\begin{itemize}[nosep]
  \item Prerequisite for economic participation (job applications, banking)
  \item Primary mode of social interaction for many
  \item Major source of information and misinformation
  \item Site of new forms of harm (cyberbullying, addiction, surveillance)
\end{itemize}

The COVID-19 pandemic revealed: digital access is now as fundamental as electricity.

\begin{intuition}[The Digital Divide]
Rural elderly: low $U^D$ (no internet, no skills).
During pandemic: couldn't access telehealth, online shopping, video calls with family.
Low $U^D$ cascaded into lower $U^P$ (health access), $U^S$ (isolation), $U^F$ (services).
\end{intuition}

\paragraph{Measurement.}
Internet access (broadband, mobile), digital literacy assessments,
e-government usage, screen time indicators, cybersecurity incident rates, data privacy indices.

%------------------------------------------------------------------------------
\subsection{Ecological Dimension ($U^{Eco}$)}
\label{sec:fepsde-ecological}

\begin{definition}[Ecological Utility]
Environmental quality and sustainability: air and water quality,
access to green spaces, climate stability, biodiversity, resource sustainability,
environmental justice.
\end{definition}

\paragraph{Why It's Increasingly Important.}
\begin{itemize}[nosep]
  \item Climate change threatens all other dimensions
  \item Environmental degradation affects health ($U^P$), livelihoods ($U^F$)
  \item Intergenerational equity: today's $U^F$ at cost of future $U^{Eco}$
  \item Nature exposure affects mental health ($U^E$)
\end{itemize}

\begin{intuition}[China's Tradeoff]
1980-2010: Rapid increase in $U^F$ (GDP growth).
Simultaneously: Severe decline in $U^{Eco}$ (air pollution, water contamination).
Health costs ($U^P$) partially offset financial gains.
Now: Policy shift toward ``ecological civilization.''
\end{intuition}

\paragraph{Measurement.}
Air quality indices (PM2.5, ozone), water quality, carbon footprint,
green space access, ecological footprint, Environmental Performance Index.

\begin{cross-reference}
See Appendix WAT (CORE-WHAT) for the complete specification of all six dimensions
including cross-dimensional complementarities and context-dependent weights.
\end{cross-reference}


%%%%%%%%%%%%%%%%%%%%%%%%%%%%%%%%%%%%%%%%%%%%%%%%%%%%%%%%%%%%%%%%%%%%%%%%%%%%%%%
%                PART III: THE THREE UTILITY CATEGORIES
%%%%%%%%%%%%%%%%%%%%%%%%%%%%%%%%%%%%%%%%%%%%%%%%%%%%%%%%%%%%%%%%%%%%%%%%%%%%%%%

\section{The Three Utility Categories: INU, KNU, IDN}
\label{sec:fepsde-categories}

%------------------------------------------------------------------------------

The six dimensions capture \emph{what} matters.
The three categories capture \emph{for whom} it matters.

\begin{definition}[Utility Categories]
\label{def:utility-categories}
\begin{align}
  INU &: \text{\textbf{Individual Utility} --- What do I want for myself?} \\
  KNU &: \text{\textbf{Collective Utility} --- What do I want for us/my group?} \\
  IDN &: \text{\textbf{Identity Utility} --- Who am I / who do I want to be?}
\end{align}
\end{definition}

%------------------------------------------------------------------------------
\subsection{Individual Utility (INU)}
\label{sec:fepsde-inu}

\begin{definition}[Individual Utility]
Standard self-interested preferences: my income, my health, my relationships.
The classical economic agent's utility function. ``What's in it for me?''
\end{definition}

\paragraph{Theoretical Grounding.}
This is the utility that Arrow-Debreu models capture.
It's real and important---but not the whole story.

%------------------------------------------------------------------------------
\subsection{Collective Utility (KNU)}
\label{sec:fepsde-knu}

\begin{definition}[Collective Utility]
Preferences over group outcomes: my family's welfare, my community's prosperity,
my nation's success, my company's performance.
\end{definition}

\paragraph{Why It's Distinct from INU.}
People sacrifice personal welfare for group welfare:
\begin{itemize}[nosep]
  \item Parents sacrifice for children
  \item Soldiers sacrifice for country
  \item Employees sacrifice for company
  \item Activists sacrifice for cause
\end{itemize}

This isn't reducible to individual utility---it's a separate category.

\begin{intuition}[Level-Indexed KNU]
The ``us'' in KNU can be indexed by level $L$:
\begin{itemize}[nosep]
  \item $KNU^1$: My household
  \item $KNU^2$: My organization/community
  \item $KNU^3$: My region/nation
  \item $KNU^\infty$: Humanity/all sentient beings
\end{itemize}
Different people weight these levels differently.
\end{intuition}

\paragraph{Theoretical Grounding.}
Connects to group selection (evolutionary theory), team reasoning (game theory),
organizational identification (management), patriotism (political science).

\begin{cross-reference}
See Appendix WHO (CORE-WHO) for the complete specification of the level hierarchy
$L \in \{0, 1, ..., \infty\}$ and the aggregation weights $\alpha^L(\Psi)$.
\end{cross-reference}

%------------------------------------------------------------------------------
\subsection{Identity Utility (IDN)}
\label{sec:fepsde-idn}

\begin{definition}[Identity Utility]
Preferences over self-concept: being a good person, being competent,
having a coherent life narrative, living according to one's values,
being respected by those one respects.
\end{definition}

\paragraph{Why It's Distinct.}
People sacrifice both INU and KNU for identity:
\begin{itemize}[nosep]
  \item Refusing to lie even when profitable and harmless
  \item Maintaining principles even when group disagrees
  \item Choosing meaningful work over higher pay
  \item Accepting hardship to ``be true to oneself''
\end{itemize}

\paragraph{Theoretical Grounding.}
Connects to:
\begin{itemize}[nosep]
  \item \citet{akerlofkranton} identity economics
  \item \citet{benaboutirole} self-signaling
  \item Self-determination theory (authenticity)
  \item Narrative identity in psychology
\end{itemize}

\begin{cross-reference}
See Appendix SCP (Social Preferences) for the formal foundations of
INU/KNU/IDN including the axioms governing category interactions.
\end{cross-reference}


%%%%%%%%%%%%%%%%%%%%%%%%%%%%%%%%%%%%%%%%%%%%%%%%%%%%%%%%%%%%%%%%%%%%%%%%%%%%%%%
%                     PART IV: TIME STRUCTURE
%%%%%%%%%%%%%%%%%%%%%%%%%%%%%%%%%%%%%%%%%%%%%%%%%%%%%%%%%%%%%%%%%%%%%%%%%%%%%%%

\section{The Time Structure: Four Horizons}
\label{sec:fepsde-time}

%------------------------------------------------------------------------------

Welfare is evaluated across time, not just at a moment.

\begin{definition}[Time Horizons]
\label{def:time-horizons}
\begin{align}
  t = 0 &: \text{\textbf{Instant} --- Right now, immediate experience} \\
  t = 1 &: \text{\textbf{Short-term} --- Days to weeks} \\
  t = 2 &: \text{\textbf{Medium-term} --- Months to years} \\
  t = 3 &: \text{\textbf{Long-term} --- Years to lifetime}
\end{align}
\end{definition}

\paragraph{Why Time Matters.}
Many choices involve intertemporal tradeoffs:
\begin{itemize}[nosep]
  \item Education: Low $U^F_0$, high $U^F_3$
  \item Exercise: Low $U^E_0$ (effort), high $U^P_2$
  \item Saving: Low $U^F_0$, high $U^F_3$
  \item Addiction: High $U^E_0$, low $U^{P,E,S,F}_2$
\end{itemize}

\begin{definition}[Discounting]
Future utility is weighted by discount factor $\delta_d(\Psi) \in (0, 1)$:
\begin{equation}
  \text{Present value of } U_{d,t} = \delta_d^t \cdot U_{d,t}
  \label{eq:discounting}
\end{equation}
\end{definition}

\begin{intuition}[Dimension-Specific Discounting]
Discount rates may vary by dimension:
\begin{itemize}[nosep]
  \item Financial: Standard exponential discounting
  \item Health: May be discounted less (``health is priceless'')
  \item Ecological: Intergenerational discounting debates
\end{itemize}
\end{intuition}


%%%%%%%%%%%%%%%%%%%%%%%%%%%%%%%%%%%%%%%%%%%%%%%%%%%%%%%%%%%%%%%%%%%%%%%%%%%%%%%
%                     PART V: GAINS AND PAINS
%%%%%%%%%%%%%%%%%%%%%%%%%%%%%%%%%%%%%%%%%%%%%%%%%%%%%%%%%%%%%%%%%%%%%%%%%%%%%%%

\section{Gains and Pains: The Prospect Theory Connection}
\label{sec:fepsde-valence}

%------------------------------------------------------------------------------

\paragraph{Why Separate Gains and Pains?}
\citet{kahnemantversky1979} showed that losses loom larger than gains:
people feel a \$100 loss more intensely than a \$100 gain.

This ``loss aversion'' applies across domains:
\begin{itemize}[nosep]
  \item Financial losses hurt more than equivalent gains help
  \item Health decline hurts more than equivalent improvement helps
  \item Social rejection hurts more than equivalent acceptance helps
\end{itemize}

\begin{definition}[Dimension-Specific Utility with Loss Aversion]
For each dimension $d$, utility is:
\begin{equation}
  U_d = G_d - \lambda_d \cdot P_d
  \label{eq:loss-aversion}
\end{equation}
where:
\begin{itemize}[nosep]
  \item $G_d \geq 0$: Gains in dimension $d$
  \item $P_d \geq 0$: Pains (losses) in dimension $d$
  \item $\lambda_d > 1$: Loss aversion coefficient for dimension $d$
\end{itemize}
\end{definition}

\paragraph{Empirical Estimates.}
\begin{itemize}[nosep]
  \item Financial: $\lambda_F \approx 2$ (Kahneman \& Tversky)
  \item Health: $\lambda_P \approx 2-3$ (health state valuations)
  \item Social: $\lambda_S \approx 2-4$ (rejection sensitivity research)
\end{itemize}

\begin{intuition}[Policy Implication]
Preventing harm is more valuable than creating equivalent benefit.
A policy that prevents \$100 loss creates more welfare than one that provides \$100 gain.
\end{intuition}


%%%%%%%%%%%%%%%%%%%%%%%%%%%%%%%%%%%%%%%%%%%%%%%%%%%%%%%%%%%%%%%%%%%%%%%%%%%%%%%
%                     PART VI: THE COMPLETE STRUCTURE
%%%%%%%%%%%%%%%%%%%%%%%%%%%%%%%%%%%%%%%%%%%%%%%%%%%%%%%%%%%%%%%%%%%%%%%%%%%%%%%

\section{The Complete 144-Component Structure}
\label{sec:fepsde-complete}

%------------------------------------------------------------------------------

\paragraph{Counting Components.}
\begin{itemize}[nosep]
  \item 3 categories (INU, KNU, IDN)
  \item 2 valences (Gain, Pain)
  \item 6 dimensions (FEPSDE)
  \item 4 time horizons
\end{itemize}

Total: $3 \times 2 \times 6 \times 4 = 144$ components.

\paragraph{Examples of Specific Components.}

\begin{table}[h]
\centering
\begin{tabular}{lll}
\toprule
\textbf{Component} & \textbf{Meaning} & \textbf{Example} \\
\midrule
$G_{INU,F,0}$ & My financial gain now & Bonus payment today \\
$P_{INU,P,2}$ & My physical pain medium-term & Chronic back pain \\
$G_{KNU,S,1}$ & Our group's social gain short-term & Team wins award \\
$P_{KNU,Eco,3}$ & Our collective ecological pain long-term & Climate change \\
$G_{IDN,E,2}$ & My identity emotional gain medium-term & Pride in achievement \\
$P_{IDN,S,0}$ & My identity social pain now & Public humiliation \\
\bottomrule
\end{tabular}
\caption{Examples of specific utility components}
\label{tab:utility-components}
\end{table}

%------------------------------------------------------------------------------
\subsection{The Aggregate Welfare Function}

\begin{definition}[Equilibrium Quality]
The quality of an equilibrium $C^*$ is:
\begin{equation}
  Q(C^*) = \sum_{i \in \{INU, KNU, IDN\}} \sum_{d \in FEPSDE} \sum_{t=0}^{3}
           \omega_i \cdot \delta_d(\Psi)^t \left[ G_{i,d,t} - \lambda_d(\Psi) \cdot P_{i,d,t} \right]
  \label{eq:equilibrium-quality}
\end{equation}
where:
\begin{itemize}[nosep]
  \item $\omega_i$: Weight on utility category $i$
  \item $\delta_d(\Psi)$: Discount factor for dimension $d$
  \item $\lambda_d(\Psi)$: Loss aversion for dimension $d$
\end{itemize}
\end{definition}

\paragraph{Parameter Estimation.}
The weights $\omega_i$, discount factors $\delta_d$, and loss aversion $\lambda_d$
can be estimated from:
\begin{itemize}[nosep]
  \item Revealed preference (observed choices)
  \item Stated preference (surveys, experiments)
  \item Deliberative methods (citizens' assemblies)
\end{itemize}

Different societies may have different parameters---this is part of context $\Psi$.

\begin{cross-reference}
See Appendix WAT (CORE-WHAT) for the complete axiom system (144 axioms)
specifying the utility structure, including Meta-Axiom U-M0 (universal utility structure).
\end{cross-reference}


%%%%%%%%%%%%%%%%%%%%%%%%%%%%%%%%%%%%%%%%%%%%%%%%%%%%%%%%%%%%%%%%%%%%%%%%%%%%%%%
%                     PART VII: INTEGRATION
%%%%%%%%%%%%%%%%%%%%%%%%%%%%%%%%%%%%%%%%%%%%%%%%%%%%%%%%%%%%%%%%%%%%%%%%%%%%%%%

\section{Integration with the CORE Framework}
\label{sec:fepsde-integration}

%------------------------------------------------------------------------------

\subsection{Connection to CORE-HOW (Complementarity)}
\label{sec:fepsde-complementarity}

\begin{cross-reference}
The complementarity parameter $\gamma$ from CORE-HOW operates on FEPSDE:
\begin{itemize}[nosep]
  \item \textbf{Inter-dimensional:} $U^F$ and $U^S$ may be complements or substitutes
  \item \textbf{Inter-category:} INU and KNU may reinforce or conflict
  \item \textbf{Cross-level:} Individual and collective utility interact via $\gamma^L$
\end{itemize}
See Chapter 5 and Appendix HOW (CORE-HOW) for the complete complementarity specification.
\end{cross-reference}

%------------------------------------------------------------------------------
\subsection{Connection to CORE-WHEN (Context)}
\label{sec:fepsde-context}

\begin{cross-reference}
The context vector $\Psi$ from CORE-WHEN modulates all FEPSDE parameters:
\begin{itemize}[nosep]
  \item Dimension weights vary by context (poverty: $U^F$ dominates)
  \item Loss aversion varies ($\lambda_d$ higher under threat)
  \item Discount rates vary (crisis: short-term focus)
\end{itemize}
See Chapter 9 and Appendix WEN (CORE-WHEN) for the 8-dimensional context specification.
\end{cross-reference}

%------------------------------------------------------------------------------
\subsection{Connection to CORE-WHO (Aggregation)}
\label{sec:fepsde-levels}

\begin{cross-reference}
The level hierarchy from CORE-WHO structures KNU:
\begin{itemize}[nosep]
  \item $L=0$: Individual (INU domain)
  \item $L=1$: Household
  \item $L=2$: Organization/Community
  \item $L=\infty$: Society/Humanity
\end{itemize}
Aggregation weights $\alpha^L(\Psi)$ determine how individual utilities combine.
See Appendix WHO (CORE-WHO) for the complete specification.
\end{cross-reference}

%------------------------------------------------------------------------------
\subsection{Connection to CORE-STAGE (BCJ)}
\label{sec:fepsde-bcj}

\begin{cross-reference}
The Behavioral Change Journey from CORE-STAGE affects FEPSDE weights:
\begin{itemize}[nosep]
  \item Stage $S_1$ (Unaware): Dimension may not enter utility at all
  \item Stage $S_4$ (Intending): Future dimensions weighted more heavily
  \item Stage $S_6$ (Maintaining): Dimension becomes automatic, low $\theta$
\end{itemize}
See Chapter 13 and Appendix STA (CORE-STAGE) for the BCJ specification.
\end{cross-reference}

%------------------------------------------------------------------------------
\subsection{The 10C Integration Table}

\begin{table}[h]
\centering
\begin{tabular}{lll}
\toprule
\textbf{CORE} & \textbf{Question} & \textbf{FEPSDE Role} \\
\midrule
WHO (AAA) & Who has utility? & Levels structure KNU \\
\textbf{WHAT (C)} & \textbf{What is utility?} & \textbf{FEPSDE dimensions} \\
HOW (B) & How interact? & Complementarity across dimensions \\
WHEN (V) & When context? & Context modulates weights \\
WHERE (BBB) & Where parameters? & 144 parameters to estimate \\
AWARE (AU) & How aware? & Awareness of each dimension \\
READY (AV) & How willing? & Willingness in each dimension \\
STAGE (AW) & Where in journey? & Stage affects dimension salience \\
\bottomrule
\end{tabular}
\caption{Integration of WHAT (FEPSDE) with the 10C framework}
\label{tab:10c-fepsde-integration}
\end{table}


%%%%%%%%%%%%%%%%%%%%%%%%%%%%%%%%%%%%%%%%%%%%%%%%%%%%%%%%%%%%%%%%%%%%%%%%%%%%%%%
%                     PART VIII: APPLICATIONS
%%%%%%%%%%%%%%%%%%%%%%%%%%%%%%%%%%%%%%%%%%%%%%%%%%%%%%%%%%%%%%%%%%%%%%%%%%%%%%%

\section{Applications and Comparisons}
\label{sec:fepsde-applications}

%------------------------------------------------------------------------------

\subsection{FEPSDE vs. Alternative Welfare Measures}

\begin{table}[h]
\centering
\begin{tabular}{llll}
\toprule
\textbf{Measure} & \textbf{Dimensions} & \textbf{Strengths} & \textbf{Weaknesses} \\
\midrule
GDP & Financial only & Widely available & Misses everything else \\
HDI & F, E, P (partial) & Simple, comparable & Only 3 dimensions \\
GNH & Multiple & Holistic & Hard to measure \\
OECD Better Life & 11 dimensions & Comprehensive & No aggregation \\
\textbf{FEPSDE} & 6 $\times$ 3 $\times$ 4 & Structured, complete & Complex \\
\bottomrule
\end{tabular}
\caption{Comparison of welfare measurement frameworks}
\label{tab:welfare-comparison}
\end{table}

\paragraph{Advantages of FEPSDE.}
\begin{enumerate}[nosep]
  \item \textbf{Structured:} Clear taxonomy enables systematic analysis
  \item \textbf{Complete:} Covers material, psychological, social, digital, ecological
  \item \textbf{Dynamic:} Explicit time structure captures sustainability
  \item \textbf{Behavioral:} Gains/pains structure incorporates Prospect Theory
  \item \textbf{Multi-level:} INU/KNU/IDN captures individual and social welfare
\end{enumerate}

%------------------------------------------------------------------------------
\subsection{Example: USA vs. Denmark}

\begin{table}[h]
\centering
\begin{tabular}{lcc}
\toprule
\textbf{Dimension} & \textbf{USA} & \textbf{Denmark} \\
\midrule
Financial ($U^F$) & Higher (GDP/capita) & Lower (after taxes) \\
Emotional ($U^E$) & Lower (life satisfaction) & Higher \\
Physical ($U^P$) & Lower (life expectancy) & Higher \\
Social ($U^S$) & Lower (trust, community) & Higher \\
Digital ($U^D$) & Similar & Similar \\
Ecological ($U^{Eco}$) & Lower (footprint) & Higher \\
\midrule
\textbf{Overall $Q$} & Depends on weights & Depends on weights \\
\bottomrule
\end{tabular}
\caption{FEPSDE comparison: USA vs. Denmark}
\label{tab:usa-denmark}
\end{table}

Which is higher depends on weights.
If $\omega_F$ dominates: USA wins.
If $\omega_{E,P,S,Eco}$ matter: Denmark wins.
FEPSDE makes the tradeoffs explicit.

%------------------------------------------------------------------------------
\subsection{The K-Q Relationship: Four Quadrants}

\begin{table}[h]
\centering
\begin{tabular}{lcc}
\toprule
& \textbf{High K (Coherent)} & \textbf{Low K (Incoherent)} \\
\midrule
\textbf{High Q (Good)} & Ideal: Nordic countries & Transition: Post-war Germany \\
\textbf{Low Q (Bad)} & Trap: North Korea & Crisis: Failed states \\
\bottomrule
\end{tabular}
\caption{The K-Q relationship: Four quadrants}
\label{tab:k-q-quadrants}
\end{table}

\begin{intuition}[The High-K, Low-Q Trap]
Being in a ``High K, Low Q'' configuration is dangerous.
The system is stable but bad.
Escaping requires accepting temporary low K (transition costs).
This explains why bad equilibria persist: the transition is costly.
\end{intuition}


%%%%%%%%%%%%%%%%%%%%%%%%%%%%%%%%%%%%%%%%%%%%%%%%%%%%%%%%%%%%%%%%%%%%%%%%%%%%%%%
%                          PART IX: SUMMARY
%%%%%%%%%%%%%%%%%%%%%%%%%%%%%%%%%%%%%%%%%%%%%%%%%%%%%%%%%%%%%%%%%%%%%%%%%%%%%%%

\section{Summary}
\label{sec:fepsde-summary}

%------------------------------------------------------------------------------

This chapter has established:

\begin{enumerate}
    \item \textbf{FEPSDE Dimensions:} Six dimensions of utility---Financial, Emotional,
          Physical, Social, Digital, Ecological---each empirically distinct and policy-relevant.

    \item \textbf{INU/KNU/IDN Categories:} Three categories capturing individual preferences,
          collective welfare, and identity concerns, with level-indexed KNU$^L$.

    \item \textbf{Time Structure:} Four horizons ($t = 0, 1, 2, 3$) with dimension-specific
          discounting capturing intertemporal tradeoffs.

    \item \textbf{Gains and Pains:} Separate tracking of gains and losses with dimension-specific
          loss aversion ($\lambda_d > 1$), incorporating Prospect Theory.

    \item \textbf{144-Component Structure:} The complete $3 \times 6 \times 4 \times 2$
          utility matrix enabling precise welfare analysis.

    \item \textbf{Quality Index $Q$:} The aggregate welfare function that ranks equilibria,
          complementing the coherence index $K$.
\end{enumerate}

%------------------------------------------------------------------------------
\paragraph{Formal Details.}
The complete formal treatment appears in Appendix~C (CORE-WHAT), which contains:
\begin{itemize}[nosep]
    \item 144 axioms specifying the utility structure
    \item Meta-Axiom U-M0 (universal utility structure)
    \item Loss aversion parameters by dimension
    \item Formal proof of completeness
\end{itemize}

%------------------------------------------------------------------------------
\paragraph{What Comes Next.}
Chapter 11 introduces \textbf{Awareness}, the cognitive dimension that determines
whether utility components enter the decision calculus at all. Without awareness,
even high potential utility in a dimension may not influence behavior.

%------------------------------------------------------------------------------
\begin{cross-reference}
\textbf{Reading Path:}
\begin{itemize}
    \item For complete FEPSDE specification: Appendix WAT (CORE-WHAT)
    \item For level aggregation: Appendix WHO (CORE-WHO)
    \item For complementarity mechanics: Chapter 5 + Appendix HOW (CORE-HOW)
    \item For context modulation: Chapter 9 + Appendix WEN (CORE-WHEN)
    \item For awareness (next): Chapter 11 + Appendix AWA (CORE-AWARE)
    \item \textbf{For intervention design: Chapter 11 (EMERGE Algorithm) + Appendix IE (EIT-EMP)}
\end{itemize}
\end{cross-reference}

%------------------------------------------------------------------------------
\paragraph{Connection to Intervention Design (EIT-EMP).}
The 144-component utility structure established in this chapter defines the \emph{target space}
for the EMERGE algorithm (Chapter 11). Each FEPSDE dimension maps to an intervention vector
$\vec{I}_d \in [0,1]$, with the INU/KNU/IDN categorization determining the \emph{emphasis}
of interventions ($I_{\text{WHAT},\text{INU}}$, $I_{\text{WHAT},\text{KNU}}$, $I_{\text{WHAT},\text{IDN}}$).
The complexity factor $f_C$ in the EMERGE formula depends directly on the number of FEPSDE
dimensions addressed by an intervention portfolio.

%%%%%%%%%%%%%%%%%%%%%%%%%%%%%%%%%%%%%%%%%%%%%%%%%%%%%%%%%%%%%%%%%%%%%%%%%%%%%%%
% END OF CHAPTER 10
%%%%%%%%%%%%%%%%%%%%%%%%%%%%%%%%%%%%%%%%%%%%%%%%%%%%%%%%%%%%%%%%%%%%%%%%%%%%%%%
